\chapter*{全书符号与约定}
\addcontentsline{toc}{chapter}{全书符号与约定}
\markboth{全书符号与约定}{全书符号与约定}

为了避免同一个对象在不同章节反复换符号，下面固定本书的主要约定。后文若无特别声明均沿用这些记号。


第一次接触群论时，真正造成阅读困难的往往不是某一个定理，而是同一行公式里同时出现了许多陌生符号。这里先约定一种读公式的方法。若写
\[
f:X\longrightarrow Y,\qquad x\longmapsto f(x),
\]
第一部分只是在说“$f$ 的输入来自集合 $X$，输出落在集合 $Y$”；第二部分才说明一个具体输入 $x$ 被送到哪里。符号 $x\in X$ 读作“$x$ 属于 $X$”，$A\subseteq B$ 读作“$A$ 是 $B$ 的子集”。若 $f$ 恰好还是群同态，$\operatorname{Ker} f$ 表示所有被 $f$ 送到单位元的输入，$\operatorname{Im} f$ 表示真正能够作为输出得到的元素。以后看到
\[
G/\operatorname{Ker} f\simeq \operatorname{Im} f
\]
时，不要把它当成一串记号，而要逐句翻译成：“把 $G$ 中被 $f$ 无法区分的元素按核的陪集合并以后，剩下的群结构与实际像完全相同。”这里 $\simeq$ 表示同构，而不是两个集合逐元素相等。

另一个容易混淆的地方是三种不同的“把空间放在一起”。集合或群的直积 $G\times H$ 的元素是有序对 $(g,h)$；向量空间的直和 $V\oplus W$ 表示把两个彼此独立的线性自由度并排放在一起，维数相加；张量积 $V\otimes W$ 则描述两个自由度同时存在，若维数有限，则
\[
\dim(V\oplus W)=\dim V+\dim W,
\qquad
\dim(V\otimes W)=\dim V\,\dim W.
\]
因此“两个能级块并列”通常对应直和，而“两个粒子、两个角动量或空间自由度与自旋自由度同时存在”通常对应张量积。后面遇到 $\oplus$ 和 $\otimes$ 时，应先问自己：现在是在“二选一地并列两个子空间”，还是在“同时给两个独立自由度各选一个状态”？

矩阵与算符记号也遵循同一个原则。$D_{ij}(g)$ 表示群元 $g$ 的表示矩阵中第 $i$ 行、第 $j$ 列的元素；$\Tr D(g)=\sum_iD_{ii}(g)$ 是迹。两个算符的交换子定义为
\[
[A,B]=AB-BA.
\]
所以 $[A,B]=0$ 的意思不是“$A$ 和 $B$ 都等于零”，而是“先做 $A$ 再做 $B$ 与反过来做得到同样结果”。量子力学中大量对称性结论都只是这句话的精确展开。

最后特别区分斜线记号的两种用法。若 $N\mathrel{\trianglelefteq} G$ 是正规子群，$G/N$ 表示带有群乘法的\emph{商群}；若 $H$ 只是一般子群，类似 $G/H$ 的写法有时只表示所有左陪集组成的\emph{集合}，未必是群。本书在第一次出现这种情形时会明确说明，以免把“陪集集合”和“商群”混为一谈。

\begin{longtable}{p{0.18\textwidth}p{0.22\textwidth}p{0.52\textwidth}}
\toprule
对象 & 主要记号 & 约定 \\
\midrule
\endfirsthead
\toprule
对象 & 主要记号 & 约定 \\
\midrule
\endhead
群与子群 & $G,H,K$ & 一般群用乘法记号；单位元统一记为 $e$。 \\
正规子群 & $N\mathrel{\trianglelefteq} G$ & 子群记为 $H\leqslant G$；正规性使用 $\trianglelefteq$。 \\
群的阶 & $|G|$ & 群元 $g$ 的阶记为 $\operatorname{ord}(g)$。 \\
共轭类 & $\operatorname{Cl}(g)$ & 中心化子记为 $C_G(g)$；中心记为 $Z(G)$。 \\
群作用 & $g\cdot x$ & 轨道与稳定子分别记为 $\operatorname{Orb}_G(x)$、$\operatorname{Stab}_G(x)$。 \\
表示 & $\rho:G\to\operatorname{GL}(V)$ & $V$ 为表示空间；具体基下矩阵写为 $D(g)$ 或 $D^{(\alpha)}(g)$。 \\
不可约表示 & $V_\alpha,\rho^{(\alpha)}$ & 维数记为 $d_\alpha$；重数记为 $m_\alpha$。 \\
特征标 & $\chi_\rho(g)$ & 定义为 $\Tr D(g)$；类函数内积用 $\langle\chi,\psi\rangle_G$。 \\
直和、张量积 & $\oplus,\otimes$ & 表示与表示空间均使用相同符号，由上下文区分。 \\
欧氏空间向量 & $\vb r,\vb a_i$ & 拉丁向量统一用直立粗体；单位向量写为 $\uv n$。 \\
点群操作 & $R$ & $R\in O(3)$；纯转动 $R\in SO(3)$。主动转动为默认约定。 \\
空间群操作 & $\{R\mid\vb t\}$ & 采用 Seitz 记号；平移子群记为 $T$。 \\
晶格、倒格子 & $\vb a_i,\vb b_i$ & $\vb a_i\cdot\vb b_j=2\pi\delta_{ij}$；倒格矢记为 $\vb G$。 \\
波矢 & $\vb k$ & $\vb k$ 与 $\vb k+\vb G$ 表示等价的 Bloch 特征标。 \\
量子对称算符 & $U(g)$ & 与抽象群元 $g$ 区分；哈密顿算符统一记为 $\hat H$。 \\
投影算符 & $P,P^{(\alpha)}$ & 满足 $P^2=P$；表示论投影另带不可约表示指标。 \\
空间转动 & $R_{\uv n}(\theta)$ & 转轴为单位向量 $\uv n$，转角为 $\theta$。 \\
旋量转动 & $U(\uv n,\theta)$ & $SU(2)$ 元素；Pauli 矩阵向量统一记为 $\boldsymbol{\sigma}=(\sigma_x,\sigma_y,\sigma_z)$。 \\
角动量 & $\vb J,\vb L,\vb S$ & 均保留 SI 量纲，显式保留 $\hbar$；量子数用 $j,l,s,m$。 \\
Wigner 矩阵 & $D^{(j)}_{m'm}$ & Euler 角约定统一为主动 $z$--$y$--$z$。 \\
置换 & $\pi,\sigma,\tau\in S_n$ & 乘法仍采用右边先作用；奇偶性记为 $\operatorname{sgn}(\pi)$。 \\
Young 分拆 & $\lambda\vdash n$ & Young 图形状写作 $\lambda=(\lambda_1,\lambda_2,\dots)$。 \\
Young 盘 & $T$ & 行群、列群记为 $R_T,C_T$；行对称子、列反对称子、Young 算符写为 $a_T,b_T,e_T$。 \\
复数单位 & $\ii$ & 自然指数统一写 $\ee^x$；不使用斜体 $i,e$ 充当这两个常数。 \\
闵可夫斯基度规 & $\eta_{\mu\nu}$ & 符号差 $(+,-,-,-)$（mostly-minus 约定）；$\mu,\nu=0,1,2,3$。 \\
洛伦兹群 & $SO^+(3,1)$ & 正时向、正取向分支；二重覆盖 $SL(2,\CC)$。 \\
洛伦兹代数 & $\mathfrak{so}(3,1)$ & 生成元 $J_i$（转动）、$K_i$（boost）；复化分解为 $A_i,B_i$。 \\
Dirac 矩阵 & $\gamma^\mu$ & Clifford 代数 $\{\gamma^\mu,\gamma^\nu\}=2\eta^{\mu\nu}$；Weyl（手征）表象为默认。 \\
手征矩阵 & $\gamma^5$ & $\gamma^5=\ii\gamma^0\gamma^1\gamma^2\gamma^3$；$\gamma^5=+1$ 左手、$-1$ 右手。 \\
Weyl 旋量 & $\psi_L,\psi_R$ & 按 $(\tfrac12,0)$、$(0,\tfrac12)$ 变换；$P_{L/R}=\tfrac12(1\mp\gamma^5)$。 \\
Poincar\'e 群 & $\mathcal P$ & $\RR^{3,1}\rtimes SO^+(3,1)$；半直积，平移为正规子群。 \\
Casimir 算符 & $P^2,W^2$ & $P^\mu P_\mu$（质量平方）、$W^\mu W_\mu$（Pauli--Lubanski）。 \\
规范群 & $G$ & 主丛结构群；规范势 $A_\mu$ 取值于 $\mathfrak g$。 \\
标准模型群 & $G_{\mathrm{SM}}$ & $SU(3)_c\times SU(2)_L\times U(1)_Y$。 \\
Dynkin 记号 & $(p,q)$ & $SU(3)$ 不可约表示标签；$p,q\in\NN$。 \\
Killing 形式 & $B(X,Y)$ & $\tr(\operatorname{ad}X\circ\operatorname{ad}Y)$；半单判据为非退化。 \\
Weyl 群 & $W$ & 根反射生成的群；$SU(N)$ 的 $W\cong S_N$。 \\
\bottomrule
\end{longtable}

本书默认在复数域上讨论有限维表示。有限群的实表示问题在需要时会明确说明；晶体学特征标表中常见的“把一对共轭一维复表示合并为二维实基组”的做法，会在附录 A 单独解释，避免把复不可约表示与实表示混为一谈。

\chapref{chap:lorentz-poincare} 起的相对论章节采用\textbf{mostly-minus 约定}：闵可夫斯基度规 $\eta_{\mu\nu}=\diag(+1,-1,-1,-1)$，四动量 $p^\mu=(E,\vb p)$，$p^2=\eta_{\mu\nu}p^\mu p^\nu=E^2-|\vb p|^2=m^2$。Dirac 方程写为 $(\ii\gamma^\mu\partial_\mu-m)\psi=0$，协变导数 $D_\mu=\partial_\mu-\ii g A_\mu^a T^a$（$T^a$ 厄米，$A_\mu^a$ 实）。此约定下 $[K_i,K_j]=-\ii\epsilon_{ijk}J_k$（boost 交换子带负号），与 \chapref{chap:rotation-group} 的 $SU(2)$ 约定 $[J_i,J_j]=\ii\epsilon_{ijk}J_k$ 一致。Lie 代数章节（\chapref{chap:lie-standard-model}）的 Killing 形式对紧群为负定，根与权的内积由 $-B$ 诱导为正定。
